% BHU PYQ -- Transcribed & Corrected
% Paper: GGRMJ-41 / GGRMJ41: Evolution of Geographical Thought
% Exam: B.A. / B.Sc. (Hons.) Semester IV Examination, 2025-26 (NEP)
% Category: Major Course
% Department: Geography

\documentclass[11pt,a4paper]{article}
\usepackage[a4paper,top=2.5cm,bottom=2.5cm,left=2.8cm,right=2.8cm,headheight=14pt]{geometry}
\usepackage{amsmath,amssymb}
\usepackage[T1]{fontenc}
\usepackage[utf8]{inputenc}
\usepackage{lmodern,microtype,fancyhdr,enumitem,hyperref,lastpage,parskip}

\hypersetup{
    colorlinks=true,
    urlcolor=black,
    linkcolor=black,
    pdftitle={GGRMJ41: Evolution of Geographical Thought -- BHU Sem IV 2025-26 (NEP)}
}

\pagestyle{fancy}
\fancyhf{}
\renewcommand{\headrulewidth}{0.4pt}
\renewcommand{\footrulewidth}{0.4pt}
\fancyhead[L]{\small\textit{Banaras Hindu University}}
\fancyhead[R]{\small\textit{GGRMJ-41: Evolution of Geographical Thought (Sem IV 2025-26)}}
\fancyfoot[C]{\small Page \thepage\ of \pageref{LastPage}}

\newlist{parts}{enumerate}{2}
\setlist[parts,1]{label=(\alph*),leftmargin=*,itemsep=4pt,topsep=2pt}
\setlist[parts,2]{label=(\roman*),leftmargin=*,itemsep=2pt,topsep=2pt}
\newcommand{\pts}[1]{\hfill\textit{[#1]}}

\begin{document}

\begin{center}
  {\large\bfseries BANARAS HINDU UNIVERSITY}\\[2pt]
  {\normalsize B.A. / B.Sc. (Hons.) --- Semester IV Examination, 2025-26 (NEP)}\\[6pt]
  \rule{0.8\linewidth}{0.5pt}\\[6pt]
  {\large\bfseries GEOGRAPHY}\\[4pt]
  {\large\bfseries Paper Code: GGRMJ-41 : Evolution of Geographical Thought}\\[4pt]
  {\normalsize Time: 2:30 Hours\hfill Maximum Marks: 70}\\[2pt]
  {\small REF. NO. 2025-26/D-IV/012}
\end{center}

\rule{\linewidth}{0.4pt}

\smallskip
\noindent\textit{Instructions: Answer \textbf{five} questions in all. All questions are of equal marks. Question No.~1 is compulsory. Answer each part of Question No.~1 in approximately 250 words and remaining questions in approximately 1000 words. Illustrate your answer with suitable maps and diagrams.}\\[4pt]
\noindent\textit{निर्देश: कुल पाँच प्रश्नों के उत्तर दीजिए। सभी प्रश्नों के समान अंक हैं। प्रश्न संख्या 1 अनिवार्य है। प्रश्न संख्या 1 के प्रत्येक भाग का उत्तर लगभग 250 शब्दों में तथा शेष प्रश्नों का उत्तर लगभग 1000 शब्दों में दीजिए। अपने उत्तर को उचित मानचित्रों एवं आरेखों द्वारा प्रदर्शित कीजिए।}

\bigskip

\noindent\textbf{Question 1.} Write short notes on the following / निम्नलिखित पर संक्षिप्त टिप्पणियाँ लिखिए:\pts{14}
\begin{parts}
  \item Geographical Contribution of Ptolemy / टॉलमी का भौगोलिक योगदान
  \item Positivism / प्रत्यक्षवाद
  \item Geographical concepts developed by Alexander Von Humboldt / अलेक्जेंडर वॉन् हम्बोल्ट द्वारा प्रतिपादित भौगोलिक संकल्पना
  \item Role of Geography in Regional Planning / प्रादेशिक नियोजन में भूगोल की भूमिका
\end{parts}

\medskip\rule{\linewidth}{0.2pt}\medskip

\noindent\textbf{Question 2.} Discuss the contribution of Arab Geographers in the development of geographical Knowledge.\\
\textit{भौगोलिक ज्ञान के विकास में अरब भूगोलवेत्ताओं के योगदान की विवेचना कीजिए।}\pts{14}

\medskip\rule{\linewidth}{0.2pt}\medskip

\noindent\textbf{Question 3.} Describe the different kinds of dichotomies existing in geographical debates.\\
\textit{भौगोलिक विमर्श में विद्यमान विभिन्न प्रकार की द्वैधताओं का वर्णन कीजिए।}\pts{14}

\medskip\rule{\linewidth}{0.2pt}\medskip

\noindent\textbf{Question 4.} Illustrate the contribution of Vidal de La Blache and French School of Geographical Thought.\\
\textit{विडाल डी ला ब्लाश और भौगोलिक चिंतन की फ्रांसीसी विचारधारा के योगदान को स्पष्ट कीजिए।}\pts{14}

\medskip\rule{\linewidth}{0.2pt}\medskip

\noindent\textbf{Question 5.} How have the great voyages contributed to the development of geographical knowledge during the 'Age of Explorations'?\\
\textit{महान यात्राओं ने 'खोजों के युग' में भौगोलिक ज्ञान के विकास में किस प्रकार योगदान दिया है?}\pts{14}

\medskip\rule{\linewidth}{0.2pt}\medskip

\noindent\textbf{Question 6.} Provide an account of the geographical knowledge that existed in ancient India.\\
\textit{प्राचीन भारत में मौजूद भौगोलिक ज्ञान का विवरण दीजिए।}\pts{14}

\medskip\rule{\linewidth}{0.2pt}\medskip

\noindent\textbf{Question 7.} Briefly describe the development of geographical thoughts after the Second World War (WW-II).\\
\textit{दूसरे विश्व युद्ध के बाद भौगोलिक चिंतन के विकास का संक्षिप्त वर्णन कीजिये।}\pts{14}

\vfill
\rule{\linewidth}{0.4pt}
\begin{center}\small
\href{https://bhu-exam.vercel.app/}{Click here to visit our website}\\[4pt]
\href{https://me-aryan.vercel.app/}{Click here to visit our contributor}
\end{center}

\end{document}
